
\chapter{Block Permutation and Separation}\label{ch:block_perm}

This chapter records the algebra-automorphism decomposition
Theorem~\ref{thm:pi_auto_decomp} that turns a permutation of the simple blocks of
a finite product algebra into a gauge transformation. The block comparison does
not order the blocks by weight modulus: equal-modulus and repeated-copy sectors
are compared through the linear-independence coefficient argument of
Chapter~\ref{ch:bnt}. The per-block linear-extension construction is stated once
in Chapter~\ref{ch:algebraic_ft}, Section~\ref{sec:pi_algebra}.

\section{Block permutation algebra}

\begin{definition}[Block ideal]\label{def:block_ideal}
    \lean{MPSTensor.blockIdeal}
    \leanok
    Let $\{R_k\}_{k=1}^r$ be simple rings.
    The $k$-th \emph{block ideal} in $\prod_{j=1}^r R_j$ is
    \[
        I_k := \{x \in \prod_{j=1}^r R_j : x_j = 0 \text{ for } j \neq k\}.
    \]
\end{definition}

\begin{theorem}[Each block ideal is an atom]\label{thm:block_ideal_atom}
    \lean{MPSTensor.isAtom_blockIdeal}
    \leanok
    \uses{def:block_ideal}
    Each block ideal $I_k$ is an atom in the lattice of two-sided ideals of
    $\prod_{j=1}^r R_j$.
\end{theorem}

\begin{proof}
    \leanok
    Under the order isomorphism between two-sided ideals of $\prod_{j=1}^r R_j$
    and products of two-sided ideal lattices, $I_k$ corresponds to the tuple with
    $\top$ in the $k$-th coordinate and $\bot$ elsewhere. This tuple is an atom,
    hence so is $I_k$.
\end{proof}

\begin{theorem}[Ring isomorphisms permute block ideals]\label{thm:ring_perm}
    \lean{MPSTensor.ringEquiv_pi_simple_permutes_blockIdeals}
    \leanok
    \uses{def:block_ideal}
    Let $\{R_k\}_{k=1}^r$ be simple rings.
    Every ring isomorphism $T : \prod_k R_k \to \prod_k R_k$
    permutes the block ideals:
    there exists a permutation $\pi$ such that
    $T$ maps the $k$-th block ideal to the $\pi(k)$-th block ideal.
\end{theorem}

\begin{proof}
    \leanok
    \uses{thm:block_ideal_atom}
    By Theorem~\ref{thm:block_ideal_atom}, each block ideal is an atom in the
    lattice of two-sided ideals of $\prod_k R_k$.
    A ring isomorphism induces an order automorphism of this lattice,
    so it permutes the atoms. Since the atoms are exactly the block ideals,
    this gives the required permutation~$\pi$.
\end{proof}

\begin{remark}[Ring level versus algebra level]\label{rem:ring_vs_alg}
    Theorem~\ref{thm:ring_perm} is a ring-level statement: it only uses the
    ideal structure of $\prod_k R_k$.
    Theorem~\ref{thm:component_map_bijective} stays at this ring level: once the
    block ideals are permuted, one can extract the corresponding single-block
    maps without using scalar multiplication.
    Theorems~\ref{thm:dim_preserved}
    and~\ref{thm:pi_auto_decomp} require an upgrade to $\C$-algebra
    automorphisms. The reason is threefold: dimension preservation compares
    the blocks as $\C$-vector spaces; Skolem--Noether applies to the resulting
    central simple $\C$-algebras; and the final per-block maps must commute
    with scalar multiplication. Thus the argument separates the ring-level
    permutation step from the later algebra-level decomposition.
\end{remark}

\begin{theorem}[Dimension preservation]\label{thm:dim_preserved}
    \lean{MPSTensor.dim_preserved}
    \leanok
    Assume $D_k \ge 1$ for all~$k$.
    If $T$ is an algebra automorphism of
    $\prod_{k=1}^r \MN{D_k}$,
    and $\pi$ is the induced permutation, then
    $D_{\pi(k)} = D_k$ for all~$k$.
\end{theorem}

\begin{proof}\leanok\uses{thm:ring_perm}
    The restriction of $T$ to the $k$-th block ideal
    induces a $\C$-algebra isomorphism
    $\MN{D_k} \cong \MN{D_{\pi(k)}}$,
    which forces $D_k = D_{\pi(k)}$
    by comparing dimensions as $\C$-vector spaces.
\end{proof}

\begin{theorem}[Per-block component maps are bijections]\label{thm:component_map_bijective}
    \lean{MPSTensor.componentMap_bijective}
    \leanok
    \uses{thm:ring_perm}
    Assume $D_k \ge 1$ for all~$k$.
    Let $T$ be a ring automorphism of $\prod_{k=1}^r \MN{D_k}$, and let $\pi$
    be the induced permutation from Theorem~\ref{thm:ring_perm}.
    For each block index $k$, the map
    \[
        M \mapsto (T(0, \ldots, 0, M, 0, \ldots, 0))_{\pi(k)}
    \]
    from $\MN{D_k}$ to $\MN{D_{\pi(k)}}$ is bijective.
\end{theorem}

\begin{proof}
    \leanok
    Theorem~\ref{thm:ring_perm} forces a single-support element in the $k$-th
    block ideal to remain supported on the $\pi(k)$-th block after applying~$T$.
    Injectivity follows from injectivity of~$T$.
    For surjectivity, apply the same argument to $T^{-1}$ and the inverse
    permutation.
\end{proof}

\begin{theorem}[Decomposition of automorphisms of $\prod \MN{D_k}$]\label{thm:pi_auto_decomp}
    \lean{MPSTensor.algEquiv_pi_matrix_decomposition}
    \leanok
    Assume $D_k \ge 1$ for all~$k$.
    Every $\C$-algebra automorphism $T$ of $\prod_{k=1}^r \MN{D_k}$
    decomposes as a block permutation~$\pi$ followed by
    per-block inner automorphisms:
    there exist $\pi \in S_r$ with $D_{\pi(k)} = D_k$
    and invertible matrices $X_k \in \GL_{D_k}(\C)$ such that,
    for every tuple $M = (M_1,\ldots,M_r)$,
    $T(M)_k = X_k M_{\pi(k)} X_k^{-1}$,
    where $T(M)_k$ denotes the $k$-th block component of $T(M)$.
\end{theorem}

\begin{proof}
    \leanok
    \uses{thm:ring_perm, thm:dim_preserved, thm:component_map_bijective,
        thm:skolem_noether}
    Theorem~\ref{thm:ring_perm} gives the permutation $\pi$.
    Theorem~\ref{thm:dim_preserved} identifies the matrix sizes
    $D_{\pi(k)} = D_k$.
    For each $k$, Theorem~\ref{thm:component_map_bijective} gives a bijective
    multiplicative map from $\MN{D_k}$ to $\MN{D_{\pi(k)}}$ extracted from the
    elements supported on one factor of~$T$; after reindexing by the dimension equality,
    this is a $\C$-algebra automorphism of $\MN{D_k}$.
    Skolem--Noether (Theorem~\ref{thm:skolem_noether}) makes each such
    automorphism inner.

    This is a standard result in the theory of semisimple algebras;
    see~\cite[Section~IV.A]{Cirac2021Matrix} for the MPS context.
\end{proof}

\section{Burnside and invariant projections}

\begin{lemma}[The algebra span reaches a finite cumulative span]\label{lem:algspan_cumulative_span}
    \lean{MPSTensor.exists_cumulativeSpan_eq_top_of_algSpan_eq_top}
    \leanok
    \uses{def:alg_span, def:cumulative_span}
    If the generated unital algebra satisfies $\algspn(A) = \MN{D}$,
    then there exists $N$ such that $T_N(A) = \MN{D}$.
\end{lemma}

\begin{proof}
    \leanok
    Every element of $\algspn(A)$ lies in some cumulative span $T_n(A)$:
    the generators lie in $T_1(A)$, scalar multiples of the identity lie in
    $T_0(A)$, and the family $\{T_n(A)\}_{n \ge 0}$ is closed under addition
    and under multiplication with lengths adding.
    Since $\MN{D}$ is finite-dimensional, the ascending chain
    $T_0(A) \subseteq T_1(A) \subseteq \cdots$ stabilizes, so one can choose a
    single $N$ with $T_N(A) = \MN{D}$.
\end{proof}

\begin{lemma}[Irreducible action gives an irreducible tensor]\label{lem:irreducible_action_to_tensor}
    \lean{MPSTensor.isIrreducibleTensor_of_isIrreducibleAction}
    \leanok
    \uses{def:irreducible_action, def:irreducible_tensor}
    If the letters of $A$ act irreducibly on $\C^D$, then $A$ has no nontrivial
    invariant orthogonal projection.
\end{lemma}

\begin{proof}
    \leanok
    If $P$ were a nontrivial invariant orthogonal projection, then its range
    would be invariant under every letter of $A$.
    Irreducibility of the action forces this range to be either $0$ or all of
    $\C^D$, so $P$ must be $0$ or $\Id$, a contradiction.
\end{proof}

\begin{theorem}[Invariant projection gives a two-block decomposition]\label{thm:lowerzero_two_block}
    \lean{MPSTensor.exists_twoBlock_decomp_of_lowerZero}
    \leanok
    \uses{def:same_mpv2}
    Suppose $P$ is an orthogonal projection satisfying, for every physical
    index~$i$,
    \[
        (\Id - P) A^i P = 0.
    \]
    Then there exist $n,m$ with $n + m = D$ and MPS tensors $A_1, A_2$ of bond
    dimensions $n$ and $m$ such that $A$ and the two-block tensor
    $A_1 \oplus A_2$ generate the same MPV family in the sense of
    Definition~\ref{def:same_mpv2}.
\end{theorem}

\begin{proof}
    \leanok
    \uses{thm:samempv_unitary}
    Diagonalize $P$ by a unitary $U$.
    Since $P^2 = P$, the diagonal form of $P$ has only $0$- and $1$-eigenvalues,
    so the basis splits as $n + m = D$.
    Conjugating $A$ by $U$ preserves the MPV family by
    Theorem~\ref{thm:samempv_unitary}.
    In this basis the relation $(\Id - P) A^i P = 0$ makes each letter upper
    block triangular.
    Discarding the off-diagonal blocks does not change the trace of any word
    product, so the resulting block-diagonal tensor has the same MPV coefficients.
    Reindex the two diagonal blocks by $\Fin n$ and $\Fin m$.
\end{proof}

\begin{theorem}[Nontrivial invariant projection gives strict two-block reduction]
    \label{thm:lowerzero_two_block_strict}
    \lean{MPSTensor.exists_twoBlock_decomp_of_lowerZero_strict}
    \leanok
    \uses{def:same_mpv2}
    Under the hypotheses of Theorem~\ref{thm:lowerzero_two_block},
    assume in addition that $P \neq 0$ and $P \neq \Id$.
    Then there exist $n,m$ with $n + m = D$, $n < D$, and $m < D$, together
    with block tensors $A_1, A_2$, such that
    $A$ and $A_1 \oplus A_2$ generate the same MPV family in the sense of
    Definition~\ref{def:same_mpv2}.
\end{theorem}

\begin{proof}
    \leanok
    \uses{thm:lowerzero_two_block}
    Apply Theorem~\ref{thm:lowerzero_two_block}.
    In the diagonal form of $P$, the $1$-eigenspace is nonempty because
    $P \neq 0$, and the $0$-eigenspace is nonempty because $P \neq \Id$.
    Hence both block sizes are positive.
    Since $n + m = D$, this implies $n < D$ and $m < D$.
\end{proof}

\section{Canonical-form predicate}

Recall from Definition~\ref{def:mpv_overlap}
and Theorem~\ref{thm:overlap_trace} that the bilinear overlap is
\[
    O_{AB}(N) = \sum_{\sigma \in \{0,\ldots,d{-}1\}^N}
    V^{(N)}(A)_\sigma \overline{V^{(N)}(B)_\sigma} = \Tr(F_{AB}^N).
\]
This notation enters the self-overlap clause of the canonical-form predicate
below.

\begin{definition}[Canonical form conditions]\label{def:is_canonical_form}
    \lean{MPSTensor.IsCanonicalForm}
    \leanok
    \uses{def:injective}
    A collection of scaling factors $(\mu_k)_{k=1}^r$ and
    block tensors $(A_k)_{k=1}^r$ satisfies the
    \emph{canonical form conditions} if:
    \begin{enumerate}
        \item each $A_k$ is injective,
        \item each $A_k$ satisfies the TP normalization
              $\sum_i (A_k^i)^\dagger A_k^i = \Id$,
        \item the moduli are non-increasing:
              $|\mu_1| \ge |\mu_2| \ge \cdots \ge |\mu_r|$,
        \item all $\mu_k \neq 0$,
        \item the self-overlap converges: $O_{A_k A_k}(N) \to 1$
              for each~$k$.
    \end{enumerate}
    Only this one-sided normalization is assumed;
    no unital condition is assumed here.
    Condition~(5) is a primitivity hypothesis:
    for a primitive channel
    (\cite[Theorem~6.7, item~3]{Wolf2012Quantum}),
    $T^n(\rho) \to \rho_\infty$ for every initial state,
    which forces the self-overlap to converge to~$1$.
    See also~\cite[Theorem~6.8]{Wolf2012Quantum}
    for the completely-positive characterisation
    and Chapter~\ref{ch:spectral} for the transfer-operator gap
    proof.
\end{definition}

\begin{remark}[Canonical-form theorem versus local predicates]\label{rem:pgvwc07_vs_local_cf}
    The translation-invariant canonical-form theorem
    of~\cite[Theorem~Th:TIcanonical, lines~742--763]{PerezGarcia2007Matrix}
    starts from an arbitrary TI-MPS representation and constructs a block
    canonical form with positive weights, unital block normalization, full-rank
    diagonal invariant densities, and uniqueness of the identity fixed point.
    Definition~\ref{def:is_canonical_form} is a local predicate used after such
    structure has already been supplied or separately proved. In particular, it
    assumes injectivity and the self-overlap limit. These results are therefore
    conditional consequences, not the source canonical-form existence or
    uniqueness theorems.
\end{remark}

\begin{remark}[Normalization convention]\label{rem:tp_vs_unital}
    The convention here differs from~\cite[Theorem~4]{PerezGarcia2007Matrix},
    which uses the unital gauge
    $\sum_i A_k^i (A_k^i)^\dagger = \Id$.
    This chapter instead works in the dual TP gauge
    $\sum_i (A_k^i)^\dagger A_k^i = \Id$.
    A positive-definite fixed point of the adjoint transfer map provides a
    gauge transformation between these two normalizations
    (Theorem~\ref{thm:left_canonical_gauge}).
\end{remark}

\begin{remark}[Overlap hypothesis]\label{rem:overlap_hypothesis}
    In the literature~\cite{Cirac2021Matrix},
    condition~(5) is derived from primitivity of the
    transfer map (peripheral spectrum $= \{1\}$).
    Theorem~\ref{thm:canonical_from_primitive} gives this implication:
    per-block primitivity implies overlap convergence,
    so condition~(5) is redundant when primitivity holds.
    The canonical form conditions retain condition~(5)
    explicitly because it can be verified independently of primitivity.
\end{remark}

\begin{remark}[Normal-canonical companion]\label{rem:normal_cf_companion}
    Definition~\ref{def:is_canonical_form} gives the block-injective conditions.
    The normal-canonical analog instead uses
    Definition~\ref{def:is_normal_canonical_form}: injectivity and the explicit
    self-overlap hypothesis are replaced by irreducibility, TP normalization,
    and primitive transfer maps. Theorem~\ref{thm:ncf_overlap_tendsto_one}
    then supplies the self-overlap limit automatically.
\end{remark}
